% problem-000054 6 碘氢体系多重平衡
\section*{6. 碘氢体系多重平衡}
\textit{下列为分问。}

\subsection*{(1)}
$\mathrm { I } _ { 2 } ( \mathrm { g } )  2 \mathrm { I } ( \mathrm { g } )$ �_{&'} = 2.00 bar (2) $\mathrm { I } _ { 2 } ( \mathrm { g } ) + \mathrm { H } _ { 2 } ( \mathrm { g } )  2 \mathrm { H I } ( \mathrm { g } )$ �_{&(} (3) HI(g) ⇌ I(g) + H(g) � = 8.0×10−6 (2) $\mathrm { H } _ { 2 } ( \mathrm { g } )  2 \mathrm { H } ( \mathrm { g } )$ �_{&\*}

\subsection*{6-1}
计算 1500 K 体系中 $\mathrm { I } _ { 2 } ( \mathrm { g } )$ 和 $\mathrm { H } _ { 2 } ( \mathrm { g } )$ 未分解时的分压。 $\left( R \mathrm { ~ = ~ } 8 . 3 1 4 \mathrm { ~ J ~ m o l ^ { - 1 } ~ K ^ { - 1 } } \right)$

\subsection*{6-2}
计算1500K平衡体系中除H(g)之外所有物种的分压。

\subsection*{6-3}
计算 $K _ { p 2 ^ { \circ } }$ 。

\subsection*{6-4}
计算 $K _ { p 4 ^ { \mathsf { c } } }$ 。（若未算出 $K _ { p 2 }$ _{(}，可设 $K _ { p 2 } = 1 0 . 0 )$
